\section{Bellman's original constructions: complete proofs}
\label{sec:annex-p17-bellman-original-constructions}

\begin{annexitem}{annex:p17.bellman.shrinking}{thm:p17-bellman-shrinking}{The exact full-domain criterion and the envelope construction}
\begin{proof}
Parameterize each allocation as \(\mathbf z=x\alpha\),
where \(\alpha\) belongs to the fixed simplex
\(\Delta=\{\alpha_i\geq0:\sum_i\alpha_i=1\}\).
This also covers \(x=0\), when every parameter gives
the same allocation \(\mathbf0\).
The simplex and the allocation domain \(D\) are compact:
a sequence of their points has a convergent subsequence
by successive bounded-coordinate subsequence extraction,
and the defining closed inequalities persist in the limit.
The elementary extraction can be obtained by nested
halves of a bounded interval in each of the finitely
many coordinates.

For continuous \(f\), the function
\(G(x,\alpha)=a(x\alpha)+f(b(x\alpha))\) is continuous.
At fixed \(x\), a maximizing sequence in \(\Delta\)
has a convergent subsequence, which attains the maximum.
If \(x_n\to x\), choose maximizing \(\alpha_n\).
Every subsequence has a further subsequence with
\(\alpha_n\to\alpha\), along which
\(Tf(x_n)=G(x_n,\alpha_n)\to G(x,\alpha)\leq Tf(x)\).
On the other hand a maximizer \(\alpha^*\) at \(x\)
gives \(Tf(x_n)\geq G(x_n,\alpha^*)\to Tf(x)\).
If \(Tf(x_n)\) failed to tend to \(Tf(x)\), a
subsequence bounded away from that value would
contradict these two observations. Thus \(Tf\)
is continuous. Also \(Tf(0)=f(0)\), since
\(a(\mathbf0)=b(\mathbf0)=0\).
Induction gives continuity and the zero boundary
for every \(F_n\). Positivity and monotonicity
of \(T\) give \(0=F_0\leq F_1\leq\cdots\).

Backward expansion of the finite maximum shows
that \(F_n(x)\) is the maximum \(n\)-stage reward.
Each finite trajectory extends indefinitely,
and nonnegative rewards make its infinite total
at least its finite total. Conversely every
partial sum on an infinite trajectory is at
most \(V(x)\). Taking suprema proves that \(V\)
is exactly the supremum infinite reward.
After \(n\) stages every remaining state is
at most \(c^n x\). Consequently, whenever
\(\eta(c^n x_0)\) is finite, every infinite
trajectory from \(x\) has reward at most
\(F_n(x)+\eta(c^n x_0)\).
Taking its supremum gives
\[
                    0\leq V(x)-F_n(x)\leq\eta(c^n x_0).
\]
If the stated small-resource condition holds,
the right side is finite for large \(n\) and
tends to zero. Thus \(V\) is finite, is the
uniform limit of continuous \(F_n\), and is
continuous with \(V(0)=0\).
The elementary maximum inequality gives
\(\|Tf-Tg\|_\infty\leq\|f-g\|_\infty\)
for bounded functions whenever these maxima
are defined. Passing to the uniform limit
in \(F_{n+1}=TF_n\) gives \(V=TV\);
its maximum is attained by the compactness
argument above.

For the converse and uniqueness, let \(F\)
be any continuous fixed point with \(F(0)=0\).
Expanding \(T^nF\) over \(n\) stages, the
only difference from \(T^n0=F_n\) is the
terminal value at a state at most \(c^n x_0\).
Hence
\[
        \|F-F_n\|_\infty
                  \leq\sup_{0\leq y\leq c^n x_0}|F(y)|
                                \longrightarrow0.
\]
The last limit is continuity at zero.
Thus \(F=V\), and its continuity at zero
gives precisely the stated condition on
\(\eta\). This also proves uniqueness,
without applying a strict sup-norm
contraction to this undiscounted operator.

For the envelope condition, at stage \(j\)
the resource is at most \(c^j x\), so its
reward is at most \(H(c^j x)\leq H(c^j x_0)\).
Every finite-horizon extension beyond stage
\(n\) has reward at most
\(\sum_{j\geq n}H(c^j x_0)\).
The assumed series therefore makes the
\(F_n\) uniformly Cauchy, with the displayed
tail bound. Its continuous limit is the
infinite value by the preceding trajectory
argument and has boundary zero.
This establishes the sufficient condition.
Here \(H(t)\to0\) as \(t\downarrow0\)
also follows directly from continuity of
\(a\) at \(\mathbf0\), since
\(\|\mathbf z\|_2\leq\sum_i z_i\leq t\).

For the linear example, let
\(K=\max_i r_i/(1-d_i)\).
Then \(r_i+Kd_i\leq K\) for every \(i\),
with equality for at least one index.
Substitution of \(F(x)=Kx\) gives
\[
 TF(x)=\max_{\sum z_i=x}\sum_i(r_i+Kd_i)z_i=Kx.
\]
The envelope is bounded by
\((\max_i r_i)t\), whose geometric sum
converges. Thus uniqueness identifies this
candidate with \(V\). Repeated allocation
to a maximizing index gives the geometric
reward sum \(r_ix/(1-d_i)=Kx\), proving
attainment constructively.
The singleton domain \(x_0=0\) has only
its prescribed zero value.
\end{proof}
\end{annexitem}

\begin{annexitem}{annex:p17.bellman.shrinking.counterexample}{thm:p17-bellman-shrinking-counterexample}{The missing between-grid control}
\begin{proof}
The \(j\)th triangular bump is supported on
\[
             [\, (5/8)4^{-j},\ (7/8)4^{-j}\,],
\]
has height \(1/(j+1)\), and is continuous
with value zero at its two endpoints.
These intervals are disjoint and accumulate
only at zero. At each positive point
only finitely many bumps meet a neighborhood.
At zero, the heights of all sufficiently
small supported bumps tend to zero.
Thus their displayed sum is a continuous
nonnegative function on \([0,1]\), zero
at zero.
Every power \(2^{-n}\) lies outside every
bump interior: the relevant neighboring
grid points at scale \(4^{-j}\) are
\((1/2)4^{-j}\) and \(4^{-j}\), whereas
the support lies strictly between them.
Thus \(a(2^{-n})=0\) for all \(n\).
For one resource there is only one allocation,
so \(h=a\) and the printed grid series is zero.

The successor \(b(x)=x/4\) is continuous
and satisfies \(b(x)\leq(1/2)x\).
From \(x=3/4\), its successive states
are the bump centers \((3/4)4^{-j}\);
their rewards are \(1/(j+1)\).
The total reward diverges: grouping terms
with indices \(2^k\) through \(2^{k+1}-1\)
gives at least \(1/2\) per group.
If a continuous fixed point with value
zero at zero existed, iterating its
one-choice equation would give
\[
 F(3/4)=\sum_{j=0}^{n-1}\frac1{j+1}
                              +F((3/4)4^{-n}).
\]
The final term tends to zero, while the
sum diverges, a contradiction.
Every original reward, continuity and
shrinking hypothesis is retained.
\end{proof}
\end{annexitem}

\begin{annexitem}{annex:p17.bellman.probability.stopping}{thm:p17-bellman-probability-stopping}{Proper policies, value iteration and uniqueness without continuity}
\begin{proof}
Fix a map \(S_l\) and its constant \(c_l\).
Consider the policy that repeatedly performs
this transformation followed by a probe,
unless absorption has already occurred.
After a transformation the nonabsorbing
probability mass is at most \(c_l\).
Thus, in each block of at most two steps,
the conditional probability of remaining
unabsorbed is at most \(c_l\).
The expected total number of steps is at
most
\[
                        2\sum_{j\geq0}c_l^j
                                      =K=\frac2{1-c_l}.
\]
If the transformation itself reaches \(e_0\),
the block stops earlier and the same bound
holds. This provides a proper policy
uniformly over the entire simplex.

For each fixed initial vector, encode the
finite action choices and probe outcomes
by a finite alphabet. A finite history
determines every state by finite composition
of the specified maps. Its probability
is a finite product of the chosen action
probabilities and probe probabilities.
Summing its possible extensions gives one,
so these laws are consistent; the countable
finite-alphabet construction, or its
real-coordinate embedding in
\cref{thm:p15-kolmogorov-countable}, defines
the path law. No measurability of a map
over all possible initial vectors is
needed for these pointwise laws.

Finite-horizon backward induction gives
\(F_n(p)=\inf_\pi E_p^\pi\min(\tau,n)\),
where \(\tau\) is absorption time.
It is attained by a deterministic
finite-horizon policy because the
action set is finite. The operator
is monotone, so \(F_n\) increases,
and the proper-policy bound gives
\(F_n\leq K\). Its pointwise limit
\(V\) is therefore finite.
Each action expression has a limit:
the probe uses only a finite sum,
and a transformation evaluates at
one fixed successor. A minimum of
finitely many convergent real
sequences tends to the minimum
of their limits. Hence \(V\)
satisfies the displayed equation.
Also \(V(e_0)=0\) and
\(V(p)\geq F_1(p)=1\) off \(e_0\).

Let \(w\) be any bounded real
solution with \(w(e_0)=0\).
It must be nonnegative. Indeed, if
\(m=\inf w<0\), every action
expression off \(e_0\) is at
least \(1+m\). Choose a negative
value \(w(p)<m+\epsilon\), with
\(0<\epsilon<\min(1,-m)\);
then \(p\ne e_0\) and
\(1+m\leq w(p)<m+\epsilon<1+m\),
a contradiction.
Thus \(w\geq0\) and \(w(p)\geq1\)
off the absorbing point.
Select a minimizing action at
each \(p\), resolving ties by
the first action in a fixed order.
Finite conditional summation
along its tree gives
\[
 w(p)=E_p\min(\tau,n)
                  +E_p[w(X_n)\mathbf1_{\{\tau>n\}}].
\]
The nonnegative remainder shows
\(E\min(\tau,n)\leq\|w\|_\infty\).
Since
\(E\min(\tau,n)=\sum_{j=0}^{n-1}P(\tau>j)\),
the infinite tail-probability sum
is finite, so absorption is almost
sure and \(E\tau\leq\|w\|_\infty\).
In particular \(P(\tau>n)\to0\).
The bounded remainder tends to
zero, proving that this greedy
policy has cost exactly \(w(p)\).

For any other policy the same
calculation gives the inequality
\[
 w(p)\leq E_p^\pi\min(\tau,n)
              +E_p^\pi[w(X_n)\mathbf1_{\{\tau>n\}}].
\]
If \(E^\pi\tau<\infty\), the
remainder tends to zero and
\(w(p)\leq E^\pi\tau\).
If its expected cost is infinite,
that comparison is automatic.
Thus every bounded solution is
the same optimal value, proving
uniqueness and greedy attainment.

Finally take an optimal \(n\)-step
policy and then switch to the
two-step-block policy above.
Its finite-horizon cost is \(F_n(p)\leq K\);
on survival past step \(n\), that
truncated cost is \(n\). Hence
\(P(\tau>n)\leq K/n\) during its
first \(n\) steps.
The expected additional cost
after switching is at most
\(K P(\tau>n)\leq K^2/n\).
Optimality of \(V\) bounds it by
this policy's total cost, proving
\(0\leq V-F_n\leq K^2/n\)
uniformly. The proof uses unit
step costs and a proper policy,
not a fictitious discount factor.
With the identity map and initial \(e_1\) in
the one-dimensional case, both actions return
to \(e_1\); the equation there becomes
\(F(e_1)=1+F(e_1)\). This proves the stated
failure when strict mass reduction is removed.
\end{proof}
\end{annexitem}

\begin{annexitem}{annex:p17.bellman.two.resources}{thm:p17-bellman-two-resources}{Local contraction and the complete interchange argument}
\begin{proof}
On any rectangle \([0,X]\times[0,Y]\),
both successor maps stay in the
rectangle. The immediate terms
\(p_1r_1x,p_2r_2y\) are nonnegative
and bounded by
\(G=p_1r_1X+p_2r_2Y\);
the continuation factors are at
most \(\bar p<1\).
The already proved
\cref{thm:p17-bellman-bounded-transformations}
therefore gives a unique bounded
fixed point on that rectangle.
Restrictions from larger rectangles
are the same fixed point on a
smaller rectangle, since successors
remain there. These compatible
solutions define the unique
locally bounded global solution.

Finite-horizon iteration from zero
is a maximum of finitely many
linear forms in \(x,y\), one for
each finite action word. It is
therefore continuous, convex and
positively homogeneous.
For any infinite word, the tail
from position \(n\) is bounded by
\(G\bar p^n/(1-\bar p)\).
This is uniform on the rectangle,
so the iterates converge uniformly
to the supremum of the infinite
returns, exactly the displayed
value \(V\). Their uniform limit
is continuous; their convexity
and homogeneity pass to the
pointwise limit.

Compare two consecutive unlike
actions at a state \((x,y)\).
The rewards for 1 then 2 are
\(p_1r_1x+p_1p_2r_2y\);
for 2 then 1 they are
\(p_2r_2y+p_2p_1r_1x\).
After both actions the state
is \((s_1x,s_2y)\), with the
same continuation multiplier
\(p_1p_2\), in either order.
Thus the first ordering is
at least as good precisely when
\[
      p_1r_1x(1-p_2)\geq p_2r_2y(1-p_1),
\]
or \(I_1\geq I_2\).
An earlier common prefix merely
multiplies this comparison by
a positive discount product.

Suppose \(I_1\geq I_2\) initially.
In any word having a first
type-1 action after \(k\) type-2
actions, move that first type-1
action left across those \(k\)
actions. Before it is used,
its resource \(x\) has not
changed; the other resource
has only decreased through
powers of \(s_2\).
Hence the same comparison
remains valid at every swap.
The return does not decrease,
and the resulting word starts
with action 1.
If a word never uses action 1,
insert one after \(k\) type-2
actions and continue with 2.
This modification changes only
a tail whose uniform bound
tends to zero as \(k\to\infty\).
Moving the inserted action to
the front gives the fixed word
1 followed forever by 2, with
at least the modified return.
Passing to the limit shows that
a word starting with 1 also
dominates the all-2 word.
Thus the supremum over all
words equals the supremum over
words starting with 1, namely
\(p_1[r_1x+V(s_1x,y)]\).
The symmetric argument proves
the second orientation. At
equality both arguments apply,
so both actions are optimal.

Along any sequence chosen by
this current-index rule, the
fixed-point equality telescopes
over each finite prefix.
The remaining value is multiplied
by at most \(\bar p^n\) and
is bounded on the invariant
initial rectangle; it tends
to zero. Hence the greedy
sequence attains \(V\).
If \(y=0<x\), the first
index is positive and the
second zero, so every greedy
action is 1. Its return is
\[
            p_1r_1x\sum_{j\geq0}(p_1s_1)^j
                             =\frac{p_1r_1x}{1-p_1s_1}.
\]
The other axis is symmetric.
At the origin the contraction
equation has value zero.
These arguments also prove
the asserted geometric error
bound and all boundary cases.
\end{proof}
\end{annexitem}

\begin{annexitem}{annex:p17.bellman.commensurate}{thm:p17-bellman-commensurate}{The finite entrance and periodic-cone formulas}
\begin{proof}
The commensurability assumption
makes \(\log s_1/\log s_2\)
positive rational. Write it
as \(u/v\) with relatively
prime positive integers and
put \(q=s_1^{1/u}=s_2^{1/v}\).
For a ratio in the stated
open cone, write \(I_1/I_2=q^{k+\theta}\),
\(0<\theta<1\).
Action 1 is optimal when the
ratio exceeds one, equivalently
\(k<0\); it multiplies the
ratio by \(q^u\), replacing
\(k\) by \(k+u\).
Otherwise action 2 multiplies
the ratio by \(q^{-v}\),
replacing \(k\) by \(k-v\).
The fractional part \(\theta\)
never changes, so the word is
the same throughout that cone.

Below \(-v\) the integer
increases by \(u\) until it
enters the interval; above
\(u-1\) it decreases by \(v\)
until it enters. The overshoot
in either case still lies in
\([-v,u-1]\).
This interval is invariant.
On its \(u+v\) integers the
map is addition of \(u\)
modulo \(u+v\), with those
integers as representatives.
Since \(\gcd(u,u+v)=1\),
it visits all representatives
once before returning: if
\(ju\) is divisible by \(u+v\),
then \(j\) is divisible by
\(u+v\), by the elementary
coprime divisibility property.
There are \(v\) negative
representatives and \(u\)
nonnegative ones, proving
the action counts.

For a finite word, direct
expansion of its successive
rewards gives \(A_wx+B_wy\).
Its remaining multiplier is
\(D_w\), and its successor
is \((s_1^{n_1(w)}x,s_2^{n_2(w)}y)\).
For one full periodic cycle
the resource multiplier is
the same in both coordinates,
\(q^{uv}\), and the discount
product is \(p_1^vp_2^u\).
The next cycle has the same
word and its rewards are
scaled by their product.
Summing this geometric series
therefore gives
\[
                 \frac{A_zx+B_zy}
                              {1-p_1^vp_2^u q^{uv}}
\]
at a cyclic starting state.
Adding the finite entrance
prefix and substituting its
successor yields exactly the
formula in the statement.
The denominator is positive
because every factor lies
strictly between zero and one.

These formulas are linear
on their cones. The powers
of \(q\) form a discrete
set of boundaries in any
compact positive ratio interval,
so the decomposition is
locally finite in the strict
positive quadrant.
At a positive boundary point,
take sequences from the two
adjacent cones. Their linear
formulas and continuity of
\(V\) give the same boundary
value. Thus no tie case is
left to an unproved choice
of formula. On the axes,
the already proved explicit
values and continuity apply;
no finite ratio is assigned.

In the numerical case,
\(u=v=1,q=1/2\).
For \(1<x/y<2\) the cycle
is 1,2, with \(A_z=1/2\),
\(B_z=1/4\), and cycle
factor \(1/8\).
Its value is \((4x+2y)/7\).
For \(1/2<x/y<1\) the
cycle is 2,1 and gives
\((2x+4y)/7\).
Continuity supplies the
endpoints of both intervals,
and substitution of \((3/2,1)\)
gives \(8/7\).
\end{proof}
\end{annexitem}

\begin{annexitem}{annex:p17.bellman.convex.split}{thm:p17-bellman-convex-split}{Convex value, every split and the exact series construction}
\begin{proof}
The source's derivative assumptions
imply the convexity used here: nonnegative
second derivative makes the
first derivative nondecreasing,
so the successive secant
slopes are ordered and the
convex chord inequality follows.
Because \(g,h\) are increasing
and vanish at zero, they are
nonnegative.
Their convex chords from zero
to \(x_0\) give
\[
       g(y)\leq y\,g(x_0)/x_0,\qquad
       h(z)\leq z\,h(x_0)/x_0.
\]
Hence \(g(y)+h(x-y)\leq Kx\).
The successor is at most
\(cx\), with \(c<1\).
The summable-envelope part of
\cref{thm:p17-bellman-shrinking}
applies to the two-component
allocation problem, giving its
unique continuous value with
boundary zero and the stated
geometric tail bound.

Start the finite-horizon
iteration at zero. If \(F_n\)
is convex, parameterize the
split by \(y=tx\),
\(0\leq t\leq1\). For fixed
\(t\), the function
\[
  g(tx)+h((1-t)x)
                 +F_n((at+b(1-t))x)
\]
is convex in \(x\).
A pointwise maximum of convex
functions is convex, by applying
their convex inequalities and
then taking the maximum.
Thus every iterate is convex.
Its uniform limit \(F\) is
convex as well.
For a fixed \(x\), the objective
as a function of \(y\) is convex:
each of \(g(y),h(x-y)\) and
\(F(bx+(a-b)y)\) is convex
under an affine substitution.
For \(0<y<x\), its convex
chord is at most the larger
of its values at \(0,x\).
Thus one of these two endpoints
attains the maximum. At \(x=0\)
there is only one split.
This proves the reduced equation
with both endpoints and all ties.

The same reasoning applies
to every finite iterate.
Thus enumerating the binary
endpoint words computes its
full \(n\)-stage value.
Every endpoint sequence leaves
at most \(c^n x\) resource
after \(n\) stages, with tail
at most \(Kxc^n/(1-c)\).
Taking the supremum of these
convergent series and passing
to the uniform finite-horizon
limit supplies the stated
exact construction.
A maximizing endpoint chosen
from \(F\) at every state
attains it: telescope its
fixed-point equalities and
use continuity of \(F\) at
the shrinking zero state.

The proof actually uses only
continuity, monotonicity,
convexity and the zero values,
so also covers \(x^\lambda\),
\(\lambda\geq1\), when its
second derivative is not
finite at zero.
Let
\(C=\max(A/(1-a^\lambda),
         B/(1-b^\lambda))\).
Then \(A+Ca^\lambda\leq C\)
and \(B+Cb^\lambda\leq C\),
with equality in at least
one of them. Thus \(Cx^\lambda\)
solves the endpoint equation
and is continuous with value
zero at zero. It is the unique
value just constructed.
Repeating an equality activity
gives its geometric power-payoff
sum. For the displayed quadratic
example the two ratios are
\(4/3\) and \(32/15\);
the latter is larger, proving
the stated value and policy.
\end{proof}
\end{annexitem}

\begin{annexitem}{annex:p17.bellman.concave.split}{thm:p17-bellman-concave-split}{The interior optimum and both closed values}
\begin{proof}
On \([0,1]\), \(g'=h'=1-x\geq0\)
and \(g''=h''=-1\).
For any split the successor
is \(x/2\), independent of
the split. The immediate
reward is
\[
 x-\frac{y^2+(x-y)^2}{2}
                   =x-\frac{x^2}{4}-(y-x/2)^2.
\]
It is uniquely maximized at
\(y=x/2\) when \(x>0\).
The reward is nonnegative
and bounded by \(x\), so
the summable-envelope theorem
gives existence and uniqueness
of a continuous value with
zero boundary.
The optimal equation is
\(F(x)=x-x^2/4+F(x/2)\).
Summing its two geometric
series gives
\[
           F(x)=\frac{x}{1-1/2}
                    -\frac{x^2/4}{1-1/4}
                         =2x-x^2/3.
\]
At either endpoint the
immediate reward would be
\(x-x^2/2\); its corresponding
sum is \(2x-2x^2/3\).
Their difference \(x^2/3\)
is positive for \(x>0\).
At zero both values and
the sole allocation are zero.
\end{proof}
\end{annexitem}

\begin{annexitem}{annex:p17.bellman.box.retrieval}{thm:p17-bellman-box-retrieval}{Ordering, primitive attempts and every posterior boundary}
\begin{proof}
For one box, the conditional
failure assumption gives
probability \(q_i^n\) of
failing on its first \(n\)
attempts, by multiplication
of the conditional probabilities.
The mean number to a completed
inspection is
\(\sum_{n\geq0}q_i^n=1/(1-q_i)\).
Its expected cost is \(c_i\).
Because this conditional law
does not depend on the ball's
location, a failure changes
no location probabilities.
A successful empty inspection
has the same likelihood \(1-q_i\)
for every other location and
zero for location \(i\).
When \(p_i<1\), normalizing
gives the stated posterior.
At \(p_i=1\) the only completed
inspection retrieves the ball,
so the continuation is zero.

For a complete inspection order
\(\sigma\), box \(\sigma(j)\)
is reached exactly when the
ball was not in an earlier
box. Its probability is
\(1-\sum_{h<j}p_{\sigma(h)}\).
Conditional on reaching it,
the expected inspection cost
is \(c_{\sigma(j)}\), so
finite summation gives the
displayed order value.
For adjacent inspected boxes
\(i,j\), let \(R\) be the
probability of reaching this
pair. Their contributions in
the two orders are
\[
            c_iR+c_j(R-p_i),\qquad
            c_jR+c_i(R-p_j).
\]
All later terms agree, since
the same two locations have
been excluded. The difference
of the first from the second
is \(p_jc_i-p_ic_j\).
Thus putting \(i\) first is
weakly better precisely when
\(p_i/c_i\geq p_j/c_j\).
Repeated swaps of inverted
adjacent pairs terminate in
a sorted order and never
increase its cost. Ties
have zero difference.
This proves the minimum
among complete orders.

Let \(F\) denote this minimum.
Partition orders by their first
box \(i\). The remaining
conditional order problem is
the same one with probabilities
normalized by \(1-p_i\).
It follows that
\[
                   F(p)=\min_i\{c_i+(1-p_i)F(p^{(-i)})\},
\]
with base \(F(e_i)=c_i\).
These equations also determine
the function uniquely by
induction on the support size.
A zero-probability box contributes
only an unnecessary positive
inspection cost and can be
placed after all positive
locations, when its probability
of being reached is zero.
Excluding it from the minimum
therefore changes no value.
All positive-support continuation
terms have smaller support;
the \(p_i=1\) term never
evaluates an undefined quotient.

It remains to prove that arbitrary
interleaving and randomization
of individual attempts cannot
do better. The collapsed equation
implies, for every allowed box,
\[
 F(p)\leq t_i+q_iF(p)
                 +(1-q_i)(1-p_i)F(p^{(-i)}).
\]
For a positive-probability box
this follows by multiplying
its collapsed inequality by
\(1-q_i\) and adding \(q_iF(p)\).
For a zero-probability box it
reduces to \(F(p)\leq t_i+F(p)\).
This is the one-attempt
inequality, including the
terminal zero when the ball
is found. Averaging it under
any randomized history policy
and telescoping through \(n\)
attempts gives
\[
 F(p)\leq E[\hbox{cost through }n\wedge\tau]
               +E[F(p_n)\mathbf1_{\{\tau>n\}}],
\]
where \(\tau\) is retrieval
time in attempts. Every
posterior value is bounded
by \(\sum_i c_i\).
If the policy has finite
expected completion cost,
positivity of \(\min_i t_i\)
implies finite expected
attempt count, hence
\(P(\tau>n)\to0\).
The remainder tends to zero
and its total expected cost
is at least \(F(p)\).
For infinite cost this
comparison is immediate.
The sorted policy repeats its
maximal index after a failure,
since no probabilities changed;
after an empty success, the
common normalization leaves
all remaining index orderings
unchanged. It completes each
visited inspection almost surely
and attains the finite order
cost already computed. Thus
it is optimal in the full
attempt-policy class.

For fixed costs and a fixed
ordering, the displayed
formula is linear in the
probabilities. Its ordering
conditions \(p_ic_j\geq p_jc_i\)
are valid throughout the
simplex. This supplies the
entire continuous prior domain,
not a grid of selected priors.
If \(q_i=1\) at a positive
prior location, every attempt
there fails surely. The ball
then cannot be retrieved on
an event of positive probability,
so expected completion time
is infinite. A prior-zero
such location remains irrelevant.

The conditional failure law is
essential to the operational
cost interpretation. If a box
is permanently blocked with
probability \(q\in(0,1)\)
and otherwise always opens,
each unrevealed marginal
failure probability can equal
\(q\), but the first failure
reveals permanent blockage.
Expected retrieval time for a
certain ball is infinite,
not \(t/(1-q)\).
The abstract completed-cost
equation and its ordering
proof remain valid when
the \(c_i\) themselves are
given; that different attempt
model does not justify
substituting the geometric
formula for them.
For the zero-cost extension, the adjacent
comparison remains \(p_jc_i-p_ic_j\).
If \(c_i=0<c_j\) and \(p_i>0\), it is
negative, so a free box should come first;
if both costs are zero the comparison is
an equality. The finite order and
completed-cost recursion therefore remain
valid without dividing by zero.
The same attempt-level verification has a
bounded remainder. For a policy required
to complete almost surely,
\(P(\tau>n)\to0\) directly, even if its
mean attempt count is infinite. Thus
the lower bound still holds, and the
sorted policy completes all its visited
geometric inspections almost surely.
This proves the extension in its declared
completion class.
\end{proof}
\end{annexitem}

\begin{annexitem}{annex:p17.bellman.box.location}{thm:p17-bellman-box-location}{Choosing the uninspected box and the highest-index alternative}
\begin{proof}
With one possible location,
the identity is already known,
so location cost is zero.
With more possible locations,
an empty completed inspection
eliminates a box, a positive
inspection ends the search,
and after all but one have
been eliminated the last
identity is inferred.
For a fixed uninspected last
box, the expected cost is
therefore the same order
sum with its final term
omitted.
Swapping two adjacent boxes
in the inspected prefix
has the same difference
\(p_jc_i-p_ic_j\) as before,
so sorting that prefix is
optimal. Taking the minimum
over the finitely many choices
of last box gives the
displayed candidate \(L\).

Equivalently, this construction
satisfies
\[
 L(p)=\min_{i:p_i>0}
          [c_i+(1-p_i)L(p^{(-i)})]
\]
when support size exceeds
one, with \(L(e_i)=0\).
The first-inspection
partition of orders proves
this recursion and induction
on the support determines
its unique value.
Multiplying the inequalities
by \(1-q_i\), adding the
failure continuation \(q_iL(p)\),
and telescoping over attempts
is exactly the preceding
retrieval verification, now
stopping when the location
is known. The remainder is
again bounded by
\((\sum_i c_i)P(\tau>n)\)
and vanishes for every
finite-cost completing policy.
Thus randomization or
interleaving failed attempts
cannot improve the candidate.
The sorted inspections attain
it almost surely.

Fix a box attaining the
largest index. In an optimal
choice of uninspected last
box, it is either that box,
or it lies in the inspected
prefix. In the latter case,
move it to the first position
by the weakly improving
adjacent swaps. This retains
optimality, proving the
source's alternative with
ties handled by existence
of such an optimal order.

For the three-box example,
the indices are \(1/4,3/10,1/20\).
The retrieval order \(2,1,3\)
has cost
\[
                1+2(7/10)+4(1/5)=16/5.
\]
Omitting box 1 gives
\(1+4(7/10)=19/5\);
omitting box 2 gives
\(2+4(1/2)=4\);
omitting box 3 gives
\(1+2(7/10)=12/5\).
This verifies the location
minimum and its policy.
In the two-box example,
the indices are \(4/5\)
and \(2/3\). Retrieval
therefore starts with box 1
and costs
\(1+(1/5)(3/10)=53/50\).
To locate the ball it is
enough to inspect the cheaper
box 2, costing \(3/10\);
if it is empty, box 1 is
inferred without inspection.
These examples exhibit both
branches of the stated
highest-index qualification.
If exactly one positive-prior box cannot be
inspected, omitting it and inspecting the
others gives a finite location procedure;
its sorted-prefix comparison is unchanged.
No completing policy gains information
by attempting that box, so it must remain
the inferred last location when necessary.
If two positive-prior boxes are uninspectable,
all possible observations have identical
laws under those two ball locations:
the inspectable boxes are empty under both,
and attempts on either uninspectable box
always fail. Their positive posterior
probabilities cannot be separated by those
observations. On the event that the ball
is in this pair, its location cannot be
identified with certainty, giving infinite
completion cost. Prior-zero boxes are
irrelevant, as before.
\end{proof}
\end{annexitem}

\begin{annexitem}{annex:p17.bellman.threshold}{thm:p17-bellman-threshold}{The minimizing integer, threshold and discontinuous boundary}
\begin{proof}
Apply the probability-vector
theorem with \(d=M=1\),
\(p=(1-x,x)\), and
\(S(1-x,x)=(1-ax,ax)\).
The probe expression is
\(1+x f(1)\) and the
transformation expression is
\(1+f(ax)\).
All its hypotheses hold
with contraction of
nonabsorbing mass by \(a\).
It gives a unique bounded
solution, nonnegative and
at least one for \(x>0\).
Write \(C=f(1)\), which
is finite and positive.

Unroll the two choices \(N\)
times before the first probe:
\[
 f(x)=\min\left\{
          n+1+a^nxC\ (0\leq n<N),\
                              N+f(a^Nx)\right\}.
\]
This follows by induction
from the equation; all
states \(a^nx\) remain
positive. The last term is
at least \(N+1\), whereas
the immediate-probe candidate
is at most \(1+C\).
For sufficiently large \(N\)
the last term cannot minimize.
Hence
\[
                   f(x)=\min_{n\geq0}(n+1+a^nxC),
\]
with a finite minimizing
index. At \(x=1\), index
zero would give \(1+C>C\),
so a minimizer \(k\) is
positive. For every \(k\geq1\),
\(C\leq k+1+a^kC\), or
\(C\leq(k+1)/(1-a^k)\);
at a minimizing index
equality holds. This proves
the stated formula for \(C\).
The ratios tend to infinity,
since each is at least
\(k+1\). A minimizer has
\(k+1\leq C\leq2/(1-a)\),
so the stated finite index
bound includes every possible
minimizer.

Let \(\tau=1/[C(1-a)]\).
The upper bound on \(C\)
gives \(\tau\geq1/2\).
For a minimizing \(k\),
\[
       1-a^k=(1-a)\sum_{j=0}^{k-1}a^j
                                      \leq k(1-a),
\]
so \(C\geq(k+1)/(k(1-a))>1/(1-a)\) and
\(\tau<1\).
Substitution of the minimizing
ratio gives the source's
displayed expression for
\(\tau\), independent of
which minimizing integer
was chosen.

For \(g_n(x)=n+1+a^nxC\),
\[
             g_{n+1}(x)-g_n(x)
                         =1-\frac{a^nx}{\tau}.
\]
These differences increase
with \(n\). They are
negative before the first
index with \(a^nx\leq\tau\),
nonnegative at that index,
and strictly positive
thereafter, except for
the one possible equality.
This proves the minimizing
index \(n(x)\) and the
stated formula for \(f(x)\).
At \(x<\tau\) immediate
probing is uniquely better;
at \(x>\tau\) at least
one transformation precedes
a minimizing probe.
At equality the neighboring
indices tie. In particular
\[
          1+\tau C=2+a\tau C
\]
because \((1-a)\tau C=1\).
Thus both displayed branches
give the same value at the
threshold, and every later
tie \(a^nx=\tau\) has the
same interpretation.
For sufficiently small
positive \(x\), \(n(x)=0\),
so \(f(x)=1+xC\to1\);
the absorbing value at
zero remains zero.

For \(a=1/2\), the first
three ratios are \(4,4,32/7\);
all later ratios are at
least five. Thus the only
minimizers are \(1,2\)
and \(\tau=1/2\).
The first shrinking step
already puts every
\(x\in(1/2,1]\) at or
below that threshold.
The two affine formulas,
their common value at
\(1/2\), and \(f(1)=4\)
follow immediately.
At \(a=0\), the equation at one gives
\(C=1+\min(C,0)\), whose only solution is
\(C=1\). It then gives \(f(x)=1\) for
all positive \(x\), with the prescribed
zero at zero. At \(a=1\), the equation
at one is \(C=1+C\), impossible for a
finite value. These endpoint calculations
explain why the source parameter range
is retained in the threshold formula.
\end{proof}
\end{annexitem}
